\scanpage{115}

\[
\deleted{T^n-1}
\]
\[
\Phi_n(T)=\prod_{\substack{\xi\ \text{d'ordre }n\\ \text{dans }\overline{\mathbb Q}}}(T-\xi)
\quad\text{polynôme \inserted{$\in\mathbb Z[T]$} unitaire de degré }\varphi(n)=\operatorname{card}(\mathbb Z/n\mathbb Z)^*
\]
\[
\begin{aligned}
\varphi(\underbrace{p_1^{\alpha_1}\cdots p_r^{\alpha_r}}_{n})
 &=\prod_i\varphi(p_i^{\alpha_i})=\prod p_i^{\alpha_i}-p_i^{\alpha_i-1}\\
 &=n\prod\left(1-\frac1{p_i}\right)
\end{aligned}
\]
\[
T^n-1=\prod_{\substack{d\in\mathbb N^*\\ d\mid n}}\Phi_d
\]
\[
\text{donc}\qquad
\Phi_n=(T^n-1)\big/\prod_{\substack{d\in\mathbb N^{\times}\\ d\mid n,\,d\ne n}}\Phi_d
\]

Si $k$ est un corps de car. \uuline{$p\nmid n$} (donc les racines de
$T^n-1=0$ y sont distinctes\ldots) alors les racines de
$\Phi_n(T)=0$ dans $k$ sont les racines \uline{primitives} $n$.ièmes
de $1$ dans $k$.

Ainsi si \inserted{\unsure{c|c'|on a}} \deleted{\unsure{on}} une réalis. \unsure{lin.}
géom. \inserted{rég.} de $\Pi_\infty$ sur $k$
d'inv. $\alpha$ tel que $\alpha=\zeta+\zeta^{-1}$,
pour qu'elle se factorise en une réalisation géom.
\uline{fidèle} de $\Pi_n$ sur $k$, \inserted{lorsque $n\geq3$,} il f. et s. que on ait
$\Phi_n(\zeta)=0$ i.e. $\Psi_n(\alpha)=0$, où $\Psi_n\in\mathbb Z[U]$ est
donné par
\[
\Psi_n(T+T^{-1})=T^{-\varphi'(n)}\Phi_n(T)
\qquad\substack{\varphi'(n)=\tfrac12\varphi(n)\\ \text{si }n\geq 3}
\]

Donc si $n$ \inserted{$\geq3$} n'est ni premier, ni double d'un \deleted{\unsure{premier}}
nb premier (donc $n=8,9,12,15,16\ldots$) alors l'anneau
modulaire pour les réalisations géom. \uline{fidèle}
\margin{$\frac{\varphi(n)}2$}
% The black disc in the interword space has no decipherable character.
% Three lines in the opposite orientation, cancelled by four strokes.
\par\noindent\rotatebox{180}{\begin{minipage}{.82\linewidth}
\deleted{Liebe Dagmar,}\\
\deleted{Vielen Dank für Deine}\\
\deleted{lieben Geburtstagsgrüsse.}
\end{minipage}}
% archive sheet number: 114
